\documentclass[10pt,a4paper]{amsart}
\usepackage[margin=19mm]{geometry}
\usepackage{amsmath,amssymb,mathtools}
\usepackage[scr=rsfs,cal=euler]{mathalfa}
\usepackage{xfrac}
\usepackage[unicode,hidelinks,bookmarks=false]{hyperref}
\newcommand{\ie}{i.e.\ }
\newcommand{\cf}{cf.\ }
\newcommand{\resp}{respectively\ }
\newcommand{\Subsection}{\subsection}
\providecommand{\nobreakdash}{\nobreak\hbox{-}}
\setlength{\emergencystretch}{2em}
\multlinegap0pt
\usepackage{bm}
\usepackage{bbm}
\usepackage{dsfont}
\usepackage{xspace}
\usepackage{cancel}
\usepackage{mathtools}
\usepackage{cleveref}
\usepackage{amsthm}
\makeatletter
\@ifundefined{theorem}{\newtheorem{theorem}{Theorem}}{}
\@ifundefined{lemma}{\newtheorem{lemma}[theorem]{Lemma}}{}
\makeatother
\newcommand{\HH}{\mathbb{H}}
\newcommand{\VV}{\mathbb{V}}
\newcommand{\la}{\langle}
\newcommand{\mI}{\mathcal{I}}
\newcommand{\ra}{\rangle}
\title[Convergence and error analysis of a semi-implicit finite vol]{Convergence and error analysis of a semi-implicit finite volume scheme for the Gray--Scott system}
\author{Tsiry Avisoa Randrianasolo}
\date{Source: arXiv:2508.18910v1 (CC BY 4.0)}
\begin{document}
\maketitle

\noindent\emph{Source and licence.} Convergence and error analysis of a semi-implicit finite volume scheme for the Gray--Scott system. Version arXiv:2508.18910v1 (CC BY 4.0), distributed under the Creative Commons Attribution 4.0 International licence, as recorded in the arXiv licence metadata for this exact version and shown on the abstract page https://arxiv.org/abs/2508.18910v1. Full rights notice: licenses/arxiv-2508.18910-CC-BY-4.0.txt.\par\smallskip

\noindent\textit{Editorial scope.} Counted unit: the Lemma whose source label is \texttt{lem:residuals} in the article (source0.tex lines 619--625, source label \texttt{lem:residuals}), delivered together with the article's own complete original proof of it (source0.tex lines 626--694, one \texttt{proof} environment of 69 lines). No mathematics is written, repaired or re-derived: statement and proof are the article's own text line for line. Required context retained uncounted: lem:approx\_pty (lines 285--291); thm:error\_estimate (lines 375--385). Every definition, display and label of those lines is retained unchanged. Omitted from this package and not redistributed: 1--284, 292--374, 386--618, 695--1077. Nothing inside a retained excerpt is omitted or altered: every retained body is delivered exactly as the article's lines are written, so no omission receipt of any kind accompanies this package. Citations: the retained excerpts carry no citation command at all, so no bibliography entry of the article is transcribed and no reference is redistributed. Reference adaptation: none was needed and none was made; every reference inside the delivered unit resolves inside it, so no unresolved reference or citation is produced. Printed slips kept literal: no printed word, symbol, hypothesis or number of the counted statement or of its proof was altered. Layout and class: the article's own class is \texttt{amsart} with packages including geometry, inputenc, babel, fontenc, lmodern, amsmath, amsfonts, amssymb, mathabx, mathrsfs, mathtools, dsfont, rotating, subcaption; the wrapper is the lane's pinned \texttt{amsart} editorial wrapper at 10pt on A4, which restates the theorem-like environments the retained text needs and adds the article's own macro definitions that the retained text needs (\texttt{\textbackslash HH}, \texttt{\textbackslash VV}, \texttt{\textbackslash la}, \texttt{\textbackslash mI}, \texttt{\textbackslash ra}), with extra wrapper packages (bm, bbm, dsfont, xspace, cancel, mathtools, cleveref). The delivered numbering is the unit's own. Assets: no retained span contains a figure environment, a graphics-inclusion command, a PDF-page inclusion, a table or a TikZ picture; no image file is copied into this package, the package declares no asset dependency and no image file is redistributed with it. Licence: version arXiv:2508.18910v1 of this article is distributed under the Creative Commons Attribution 4.0 International licence, as recorded in the arXiv licence metadata for this exact version and shown on the abstract page https://arxiv.org/abs/2508.18910v1. A full rights notice is delivered at licenses/arxiv-2508.18910-CC-BY-4.0.txt. Characters: every retained excerpt is pure ASCII, so the delivered engine, XeLaTeX with its default Latin Modern text font, reports no missing character.
	\begin{lemma}\label{lem:approx_pty}
		For any $u\in H^1$, the following inequality holds
		\begin{equation*}
			\Vert u-\mI_h u\Vert_{L^2}^2\leq \gamma_0h^2\Vert\nabla u\Vert^2_{L^2},
		\end{equation*}
		where the constant $\gamma_0$ depends only on the reference cell geometry and the quadrature rule.
	\end{lemma}

	\begin{theorem}[Error estimate]\label{thm:error_estimate}
		With further regularity, i.e.,
		\[
		u, v \in L^\infty(0,T; H^2\cap\HH),\quad \partial_t u, \partial_t v \in L^2(0,T; \VV),\quad \partial_{tt} u, \partial_{tt} v \in L^\infty(0,T; \VV');
		\]
		we have
		\begin{equation*}
			\max_{0\leq n\leq N} \Big(\Vert u_h^n - u(t^n) \Vert_{\HH} + \Vert v_h^n - v(t^n) \Vert_{\HH}\Big)\leq C(h^2 + \Delta t),
		\end{equation*}
		where the constant $C>0$ depends on the domain $D$, terminal time $T$, and norms of the exact solution $u, v$, but is independent of $h$, $\Delta t$, or $n$.
	\end{theorem}

	\begin{lemma}\label{lem:residuals}
		Assume that the hypotheses of \cref{thm:error_estimate} hold. Then, the residuals satisfy  
		\begin{equation*}
			\max_{0\leq n\leq N} \Big(\Vert R_h^n(\phi_h)\Vert_{H^{-1}} + \Vert Q_h^n(\psi_h)\Vert_{H^{-1}}\Big)\leq C(h^2 + \Delta t),\quad\forall \phi_h,\psi_h\in \VV_h,
		\end{equation*}
		where $C>0$ is a constant that depends on the norms of $u,v$, but independent on the scheme.
	\end{lemma}

	\begin{proof}
		We give the proof for $R_h^n$; the case of $Q_h^n$ follows analogously.
		Let $\phi_h\in \VV$. To estimate the action of $R_h^n$ on $\phi_h$, we insert the weak form and subtract it, then bound the dual norm. We split the residual into 3 terms,
		\begin{equation*}
			R_h^n(\phi_h) = T_1 + T_2 + T_3,
		\end{equation*}
		with  
		\begin{align*}
			\intertext{the \textit{time truncation}}
			T_1&\coloneqq \Big( \frac{u(t^{n+1}) - u(t^n)}{\Delta t} - \partial_t u(t^{n+1}),\, \phi_h \Big)_h  +  \big(\partial_t u(t^{n+1}),\, \phi_h\big)_h - \big\la\partial_t u(t^{n+1}),\, \phi_h \big\ra,
			\\
			\intertext{  the \textit{gradient quadrature error}}
			T_2&\coloneqq d_u \big( \nabla u(t^{n+1}), \,\nabla \phi_h \big)_h- d_u \big( \nabla u(t^{n+1}), \,\nabla \phi_h \big),
			\vphantom{\Big( \frac{u(t^{n+1}) - u(t^n)}{\Delta t} - \partial_t u(t^{n+1}),\, \phi_h \Big)_h}
			\\
			\intertext{and the \textit{nonlinear quadrature error}}
			T_3&\coloneqq \big( f(u(t^n),v(t^n)), \, \phi_h \big)_h - \big(f(u(t^n),v(t^n)),\, \phi_h \big).
			\vphantom{\Big( \frac{u(t^{n+1}) - u(t^n)}{\Delta t} - \partial_t u(t^{n+1}), \,\phi_h \Big)_h}
		\end{align*}
		

		\textit{Time truncation.} 
		Using Taylor expansion,
		\begin{equation*}
			\frac{u(t^{n+1}) - u(t^n)}{\Delta t} - \partial_t u(t^{n+1})  = \frac{\Delta t}{2}\partial_{tt}u(\xi),\quad\mbox{ for some } \xi\in (t^{n}, t^{n+1}),
		\end{equation*}
		which gives for the first term
		\begin{equation*}
			\bigg\Vert \frac{u(t^{n+1}) - u(t^n)}{\Delta t} - \partial_t u(t^{n+1}) \bigg\Vert_{H^{-1}}\leq C\Vert \partial_{tt}u\Vert_{L^\infty(t^n,t^{n+1}; H^{-1})}.
		\end{equation*}
		We use \cref{lem:approx_pty} on the second term (quadrature error for $L^2$-inner product), which gives
		\begin{equation*}
			\big\vert  \big(\partial_t u(t^{n+1}),\, \phi_h\big)_h - \big\la\partial_t u(t^{n+1}),\, \phi_h \big\ra\big\vert \leq Ch^2\big\Vert \partial_t u(t^{n+1})\big\Vert_{H^1}\Vert \phi_h\Vert_{H^1}.
		\end{equation*}
		Thus, the first contribution is
		\begin{equation*}
			\vert T_1\vert\leq C(\Delta t + h^2) \Vert \phi_h\Vert_{H^1}.
		\end{equation*}
		
		\textit{Gradient quadrature error.} Since $\nabla u(t^{n+1})\in H^1\times H^1$, and  we use midpoint rule (or piecewise constant projection), by \cref{lem:approx_pty}, the quadrature error satisfies
		\begin{equation*}
			\big\vert   \big( \nabla u(t^{n+1}), \nabla \phi_h \big)_h-   \big( \nabla u(t^{n+1}), \nabla \phi_h \big)\big\vert \leq C h^2\Vert \nabla u(t^{n+1})\Vert_{H^{1}}\Vert \phi_h\Vert_{H^1}.
		\end{equation*}
		Thus, 
		\begin{equation*}
			\vert T_2\vert\leq Ch^2\Vert \phi_h\Vert_{H^{1}}.
		\end{equation*}
		
		\textit{Nonlinear quadrature error.} The reaction term $f(u,v) = -uv^2 + F(1-u)\in H^1$ if $u,v\in H^2$. Similarly, by applying \cref{lem:approx_pty}, we have
		\begin{equation*}
			\big\vert  \big(f(u(t^n),v(t^n)), \, \phi_h \big)_h - \big(f(u(t^n),v(t^n), \, \phi_h\big)\big\vert \leq Ch^2\Vert f(u,v)\Vert_{H^1}\Vert \phi_h\Vert_{H^{1}}.
		\end{equation*}
		Thus, 
		\begin{equation*}
			\vert T_3\vert \leq Ch^2\Vert \phi_h\Vert_{H^{1}}.
		\end{equation*}
		
		Combining the estimates, we arrive at
		\begin{equation*}
			\vert R_h^n(\phi_h)\vert \leq C(\Delta t + h^2)\Vert \phi_h\Vert_{H^{1}},
		\end{equation*}
		thus, by definition of the dual norm
		\begin{equation*}
			\Vert R_h^n(\phi_h)\Vert_{H^{-1}}\leq C(h^2 + \Delta t),
		\end{equation*}
		where $C>0$ is a constant depends on the norms of $u$, but not on the scheme.
		
		That completes the proof of the lemma.
	\end{proof}

\end{document}
